\section{向量空间\texorpdfstring{$\R^{n}$}{Rn}}

\begin{definition}{}{}
	定义$\R$在$\R^{n}$上的左作用$\cdot_{1}$和右作用$\cdot_{2}$如下:对任一$t\in\R$,$\left(x_{1},\ldots,x_{n}\right)\in\R^{n}$,
	\begin{align*}
		t\cdot_{1}\left(x_{1},\ldots,x_{n}\right)=\left(tx_{1},\ldots,tx_{n}\right),\left(x_{1},\ldots,x_{n}\right)\cdot_{2}t=\left(x_{1}t,\ldots,x_{n}t\right),
	\end{align*}
	则$\R^{n}$是$\left(\R,\R\right)$-双模.此外,$\R^{n}$按$\cdot_{1}$成为拓扑左$\R$-模,按$\cdot_{2}$成为拓扑右$\R$-模.今后我们将不再区分$\cdot_{1}$和$\cdot_{2}$.
\end{definition}

\begin{definition}{}{}
	对任一$i\in\llbracket 1,n\rrbracket$,记$e_{i}=\left(\delta_{1,i},\ldots,\delta_{i,i},\ldots,\delta_{n,i}\right)$.我们称$\left(e_{i}\right)_{i=1}^{n}$是$\R^{n}$的标准有序基.
\end{definition}

\begin{proposition}{}{}
	$\R^{n}$是$n$维$\R$-向量空间.
\end{proposition}

\begin{proposition}{}{}
	设$m\in\Z_{\geqslant 1}$,$f:\R^{n}\to\R^{m}$是仿射映射,则$f$按$\R^{n}$和$\R^{m}$的加法一致结构一致连续.
\end{proposition}

\begin{corollary}{}{}
	设$f:\R^{n}\to\R^{n}$是仿射同构,则$f$是同胚,也是一致同构.
\end{corollary}

\begin{definition}{}{}
	设$\left(a_{i}\right)_{i=1}^{n}$是$\R^{n}$的有序基,$b\in\R^{n}$.记
	\begin{align*}
		P_{1}=\left\{b+\sum\limits_{i=1}^{n}u_{i}a_{i}\in\R^{n}\middle|\forall i\in\llbracket 1,n\rrbracket\left(-1\leqslant u_{i}\leqslant 1\right)\right\},P_{2}=\left\{b+\sum\limits_{i=1}^{n}u_{i}a_{i}\in\R^{n}\middle|\forall i\in\llbracket 1,n\rrbracket\left(-1<u_{i}<1\right)\right\}.
	\end{align*}
	依次称$P_{1}$,$P_{2}$是以$b$为中心,$\left(a_{i}\right)_{i=1}^{n}$为基的闭平行多面体和开平行多面体.显然$P_{1}$是$b$的紧邻域.
\end{definition}